\documentclass{article}
\usepackage{amsmath}
\usepackage{siunitx}

\title{High-Frequency Decoupling: The Mathematics of Parallel Capacitance}
\author{}
\date{}

\begin{document}
\maketitle
\section{Capacitors suck at doing what they should}
A capacitor does not actually store energy perfectly like we are taught. It's actually an LCR circuit with equivalent series resistance ($R_{ESR}$) and equivalent series inductance ($L_{ESL}$).\\

The complex impedance $Z$ as a function of angular frequency $\omega$ is as such:
\begin{equation}
    Z(\omega) = R_{ESR} + j\omega L_{ESL} + \frac{1}{j\omega C}
\end{equation}

The size of this impedance, which is how effective the component is at providing current during a voltage drop, is:
\begin{equation}
    |Z| = \sqrt{R_{ESR}^2 + \left(\omega L_{ESL} - \frac{1}{\omega C}\right)^2}
\end{equation}

\section{Frequency}
Because inductive reactance ($\omega L_{ESL}$) is proportional to frequency while capacitive reactance ($\frac{1}{\omega C}$) is inversely proportional, they should perfectly cancel each other out at a frequency called the self resonant frequency ($f_{SRF}$):
\begin{equation}
    f_{SRF} = \frac{1}{2\pi\sqrt{L_{ESL} C}}
\end{equation}

If $f > f_{SRF}$, the $j\omega L_{ESL}$ wins. The component stops acting as a storage device and more like an inductor opposing change in current. A \SI{10}{\micro\farad} capacitor is big, resulting in a high $L_{ESL}$ and a low $f_{SRF}$.

\section{Logic}
A \SI{2}{\mega\hertz} SPI clock is a square wave rather than the sin waves most are accustomed to. Fourier analysis says that a square wave is really just an infinite sum of odd sine harmonics:
\begin{equation}
    v(t) = \frac{4V_{DD}}{\pi} \sum_{n=1,3,5\dots}^{\infty} \frac{\sin(n \omega t)}{n}
\end{equation}


The sharp voltage edges demand instantaneous current according to:
\begin{equation}
    i(t) = C \frac{dV}{dt}
\end{equation}
Simply put, because change in voltage in square waves are nearly instant, the current is super high.

Because the higher order harmonics generated by these jumps (quick change to voltage) easily exceed the low $f_{SRF}$ of the \SI{10}{\micro\farad} capacitor, the inductance actually kills our current. This causes the voltage rail to fail according to Faraday's law of induction:
\begin{equation}
    \Delta V = L_{ESL} \frac{di}{dt}
\end{equation}

\section{Why we need 2 capacitors}
By placing $C_1$ (\SI{10}{\micro\farad}) and $C_2$ (\SI{0.1}{\micro\farad}) in parallel across the 3v3 to GND line, we control the equivalent impedance $Z_{eq}$ using total admittance $Y_{eq}$:
\begin{align}
    Y_{eq}(\omega) &= \frac{1}{Z_1(\omega)} + \frac{1}{Z_2(\omega)} \\
    Z_{eq}(\omega) &= \frac{Z_1(\omega) Z_2(\omega)}{Z_1(\omega) + Z_2(\omega)}
\end{align}

At low frequencies, $C_1$ has massive charge storage, which has a low $Z_{eq}$. At high frequencies ($\omega > \omega_{SRF1}$), $Z_1$ becomes inductive and blocks current. However, because $C_2$ is microscopic in comparison, its $L_{ESL}$ is near 0, pushing its $f_{SRF}$ into the high Megahertz range.

$Z_2$ remains strictly capacitive at these upper frequencies, dominating the parallel equation and forcing $Z_{eq}(\omega)$ to remain low enough to ignore the high speed demands.

\end{document}
